\clearpage
\phantomsection
\section*{Chương 3. DC-blocker IIR (cực/zero, accumulator 32-bit)}
\addcontentsline{toc}{section}{Chương 3. DC-blocker IIR (cực/zero, accumulator 32-bit)}
\pdfbookmark[1]{Chương 3. DC-blocker IIR (cực/zero, accumulator 32-bit)}{chuong3}

\begin{learningobj}
Sau Chương~3, người học giải thích được vì sao offset một chiều phá đỉnh tự tương quan, viết đúng hàm truyền một cực--một zero với $\alpha=1-2^{-8}$, hiện thực RTL accumulator 32~bit theo công thức $y_{\mathrm{acc}}\leftarrow y_{\mathrm{acc}}+x_{\Delta}-(y_{\mathrm{acc}}\ggg 8)$ rồi $\mathrm{sat}_{16}(y_{\mathrm{acc}}\ggg 15)$, và đối chiếu bit-exact với twin Python. Khái niệm điểm cực/điểm không và Avalon-ST đã có ở ô Feynman~12 của Chương~0.
\end{learningobj}

\subsection*{3.1. Hiện tượng: offset DC nuốt đỉnh chu kỳ}

\begin{physicalinsight}
Micro, codec WM8731 và đường dây analog luôn cộng thêm một thành phần gần đứng yên: điện áp offset một chiều (DC), cộng thêm trôi chậm do nhiệt. Trong miền mẫu, đó là một hằng $A$ cộng vào mọi $x[n]$. Phép tự tương quan (ô Feynman~5, Chương~0) không phân biệt năng lượng hữu ích và năng lượng offset. Nếu $A$ lớn, đỉnh tại độ trễ $\tau=0$ phình theo $A^2$ và có thể át hoàn toàn đỉnh chu kỳ $T_0$ của dây thanh --- thuật toán sẽ báo vô thanh hoặc treo ở cao độ sai.
\end{physicalinsight}

Viết tường minh với tín hiệu $x[n]+A$:
\begin{align*}
\mathcal{R}_{x+A}[0]
&=\sum_{n=0}^{N-1}(x[n]+A)^2\\
&=\sum_{n=0}^{N-1}x[n]^2 + 2A\sum_{n=0}^{N-1}x[n] + N A^2
=\mathcal{R}_x[0]+2A\sum_n x[n]+N A^2.
\end{align*}
Số hạng $N A^2$ luôn dương và tỉ lệ với bình phương offset. Praat trừ mean cục bộ từng khung; trên FPGA streaming ta dùng bộ lọc IIR liên tục gọi là \textbf{DC-blocker}: một bộ thông cao rất nhẹ, ổn định, rẻ silicon, đặt ngay sau I\textsuperscript{2}S/CDC.

\begin{figure}[H]
\centering
\begin{tikzpicture}[font=\small]
  \begin{scope}
    \draw[->] (0,0) -- (5.2,0) node[right] {$\tau$};
    \draw[->] (0,0) -- (0,2.8) node[above] {$\mathcal{R}$};
    \draw[thick, red!70!black] (0,2.4) -- (0.15,2.4) -- (0.3,0.4) -- (1,0.35) -- (2,0.9) -- (3,0.3) -- (4,0.25);
    \node[font=\footnotesize\bfseries] at (2.5,3.15) {Trước DC-block: $\mathcal{R}[0]$ phình};
  \end{scope}
  \begin{scope}[xshift=7.2cm]
    \draw[->] (0,0) -- (5.2,0) node[right] {$\tau$};
    \draw[->] (0,0) -- (0,2.8) node[above] {$\mathcal{R}$};
    \draw[thick, teal!70!black] (0,1.2) -- (0.2,0.3) -- (1,0.25) -- (2,1.8) -- (3,0.3) -- (4,0.25);
    \node[font=\footnotesize\bfseries] at (2.5,3.15) {Sau: đỉnh $T_0$ lộ rõ};
    \node[teal!60!black, font=\scriptsize] at (2,2.15) {$T_0$};
  \end{scope}
\end{tikzpicture}
\caption{Cùng một tín hiệu có offset: bên trái đỉnh $\tau=0$ nuốt chu kỳ; bên phải sau DC-blocker đỉnh $T_0$ hiện ra để Chương~9 nội suy.}
\label{fig:ch03_dc_before_after}
\end{figure}

\paragraph{Ví dụ tính tay --- DC thuần phá $\mathcal{R}[0]$.} Chuỗi $x=[1,1,1,1]$ (toàn DC) cho $\mathcal{R}[0]=\sum x_n^2=4$. Sau khi trừ mean (mọi mẫu về 0) thì $\mathcal{R}[0]=0$: DC thuần phải đi tiếp như tín hiệu vô thanh, không được nuôi FSM cao độ bằng năng lượng giả.

\subsection*{3.2. Toán: cực, zero và chọn $\alpha=1-2^{-8}$}

Hàm truyền chuẩn của DC-blocker bậc nhất:
\[
H(z)=\frac{1-z^{-1}}{1-\alpha z^{-1}},\qquad 0<\alpha<1.
\]
\textbf{Zero} (nghiệm tử số) tại $z=1$ --- đúng tần số $0\,\text{Hz}$ --- nên $H(1)=0$: mọi tín hiệu hoàn toàn đứng yên bị triệt tiêu. \textbf{Cực} (nghiệm mẫu số) tại $z=\alpha$, nằm sát zero nhưng vẫn trong vòng đơn vị nên bộ lọc ổn định. Phương trình sai phân tương đương:
\[
y[n]=x[n]-x[n-1]+\alpha\, y[n-1].
\]

\begin{figure}[H]
\centering
\begin{tikzpicture}[scale=1.35]
  \draw[thick] (0,0) circle (1.5);
  \draw[->] (-2.0,0) -- (2.2,0) node[right] {$\mathrm{Re}$};
  \draw[->] (0,-1.9) -- (0,1.9) node[above] {$\mathrm{Im}$};
  \node[below left] at (0,0) {$0$};
  \draw[very thick, blue!70!black] (1.5,0.12) -- (1.5,-0.12);
  \draw[very thick, blue!70!black] (1.38,0) -- (1.62,0);
  \node[above right,blue!70!black] at (1.5,0.15) {zero $z=1$};
  \filldraw[red!70!black] (1.494,0) circle (0.055);
  \node[below,red!70!black,font=\small] at (1.35,-0.18) {cực $z=\alpha=1-2^{-8}$};
  \node[align=left, font=\small] at (-0.15,-2.35)
    {Zero trên vòng đơn vị tại DC $\Rightarrow H(1)=0$.\\
     Cực gần zero (trong vòng) $\Rightarrow$ HPF nhẹ, ổn định.};
\end{tikzpicture}
\caption{DC-blocker trên mặt phẳng $z$: zero tại $z=1$, cực tại $\alpha=1-2^{-8}=0{,}99609375$.}
\label{fig:ch03_zplane}
\end{figure}

Trên silicon ta chọn $\alpha=1-2^{-8}=0{,}99609375$ vì phép nhân $\alpha\cdot y$ thu gọn thành một phép dịch và trừ:
\[
\alpha y = y - (y\ggg 8) = y\cdot(1-2^{-8}),
\]
không tốn khối DSP. Hằng thời gian xấp xỉ $1/(1-\alpha)=256$ mẫu. Tại $f_s=16\,\text{kHz}$ (sau FIR), cutoff xấp xỉ
\[
f_c\approx\frac{1-\alpha}{2\pi}f_s=\frac{2^{-8}}{2\pi}\cdot 16000\approx 9{,}95\,\text{Hz},
\]
nằm dưới $F_{0,\min}=75\,\text{Hz}$ và trên DC --- đúng chỗ cần chặn. (Ở $48\,\text{kHz}$ ngay sau codec, $f_c\approx 29{,}8\,\text{Hz}$, vẫn an toàn trước khi hạ mẫu.)

\paragraph{Ví dụ tính tay --- đáp ứng xung với $\alpha=0{,}5$ (cực xa để thấy nhanh).} Cho $x=\delta[n]$ hiệu dụng qua công thức $y[n]=(x[n]-x[n-1])+\alpha y[n-1]$ với đầu vào hằng $x=[1,1,1,1,1]$, $y[-1]=0$:
\begin{center}\small
\begin{tabular}{@{}rlll@{}}
\toprule
$n$ & $x[n]$ & $y[n]=(x[n]-x[n-1])+\alpha y[n-1]$ & giá trị \\
\midrule
0 & 1 & $1-0+0$ & $1$ \\
1 & 1 & $0+0{,}5\cdot 1$ & $0{,}5$ \\
2 & 1 & $0+0{,}5\cdot 0{,}5$ & $0{,}25$ \\
3 & 1 & $0+0{,}5\cdot 0{,}25$ & $0{,}125$ \\
4 & 1 & $0+0{,}5\cdot 0{,}125$ & $0{,}0625$ \\
\bottomrule
\end{tabular}
\end{center}
Năng lượng bị nuốt theo cấp số nhân về 0 --- đúng trực giác HPF nhớ trạng thái. Với $\alpha=0{,}99609375$ suy giảm chậm hơn khoảng $256$ lần mẫu.

\subsection*{3.3. Phần cứng: vì sao accumulator 32-bit và RTL}

FIR Compiler II không thay thế được IIR có phản hồi: kiến trúc FIR không có vòng $y[n-1]$. Viết khoảng mười dòng Verilog với accumulator 32~bit là lựa chọn dạy và lab.

Nếu giữ $y$ chỉ 16~bit Q1.15, mỗi lần nhân $\alpha$ (hay dịch $\ggg 8$) sẽ cắt mất phần thập phân khi $|y|$ nhỏ, tạo vùng chết (deadband) hoặc dao động limit-cycle. Giải pháp: giữ thanh ghi phản hồi \texttt{y\_acc} signed 32~bit; mẫu vào Q1.15 được kéo lên cùng thang bằng dịch trái 15~bit trước khi trừ; ngõ ra lấy lại Q1.15 bằng dịch phải 15~bit rồi bão hòa.

\paragraph{Phương trình bit-exact (khớp assignment và twin).}
\begin{align*}
x_{\Delta} &= (x[n]\ll 15)-(x[n-1]\ll 15),\\
y_{\mathrm{acc}}[n] &= y_{\mathrm{acc}}[n-1] + x_{\Delta} - (y_{\mathrm{acc}}[n-1]\ggg 8),\\
y_{\mathrm{out}} &= \mathrm{sat}_{16}(y_{\mathrm{acc}}[n]\ggg 15),
\end{align*}
với $\mathrm{sat}_{16}$ ghim vào $[-32768,+32767]$. Trên số có dấu, toán tử Python/NumPy \texttt{>>} và Verilog \texttt{>>>} cùng là arithmetic floor shift.

\begin{verbatim}
module dc_blocker (
  input  wire               clk, rst_n,
  input  wire               in_valid,
  output wire               in_ready,
  input  wire signed [15:0] x_in,       // Q1.15
  output reg                out_valid,
  input  wire               out_ready,
  output reg  signed [15:0] y_out
);
  reg signed [15:0] x_z1;
  reg signed [31:0] y_acc;              // 32-bit feedback (NOT 16)
  assign in_ready = out_ready | ~out_valid;

  wire signed [31:0] x_delta =
      ($signed({x_in, 15'b0})) - ($signed({x_z1, 15'b0}));
  wire signed [31:0] y_next =
      y_acc + x_delta - (y_acc >>> 8);
  wire signed [31:0] y_shift = y_next >>> 15;

  function automatic signed [15:0] sat16;
    input signed [31:0] v;
    begin
      if (v > 32'sd32767)       sat16 = 16'sd32767;
      else if (v < -32'sd32768) sat16 = -16'sd32768;
      else                      sat16 = v[15:0];
    end
  endfunction

  always @(posedge clk or negedge rst_n) begin
    if (!rst_n) begin
      x_z1 <= 0; y_acc <= 0; y_out <= 0; out_valid <= 0;
    end else if (in_valid && in_ready) begin
      y_acc <= y_next;
      x_z1  <= x_in;
      y_out <= sat16(y_shift);
      out_valid <= 1'b1;
    end else if (out_valid && out_ready) out_valid <= 1'b0;
  end
endmodule
\end{verbatim}

\subsection*{3.4. Twin Python bit-exact và ví dụ step Q1.15}

\begin{verbatim}
import numpy as np

def sat16(v):
    return int(np.clip(v, -32768, 32767))

def dc_blocker_q15(x, alpha_shift=8):
    """x: int16 Q1.15. y_acc int32.
    y_acc += x_delta - (y_acc >> 8); out = sat16(y_acc >> 15)
    với x_delta = (x<<15) - (x_z1<<15)."""
    y_acc = np.int32(0)
    x_z1 = np.int32(0)
    out = np.empty(len(x), dtype=np.int16)
    for i, xi in enumerate(np.asarray(x, dtype=np.int32)):
        dx = np.int32((xi << 15) - (x_z1 << 15))
        y_acc = np.int32(y_acc + dx - (y_acc >> alpha_shift))
        x_z1 = xi
        out[i] = sat16(int(y_acc >> 15))
    return out
\end{verbatim}

\paragraph{Ví dụ tính tay khớp bit.} Đầu vào $[0,0,0]$ rồi bước DC $16384$ ($0{,}5$ trong Q1.15). Ba mẫu đầu ra $0$. Tại mẫu bước ($n=3$): $x_{\Delta}=16384\ll 15=536870912$, $y_{\mathrm{acc}}=536870912$, $y_{\mathrm{out}}=16384$. Mẫu kế:
\[
y_{\mathrm{acc}}\leftarrow 536870912 + 0 - (536870912\ggg 8)=534773760,\quad
y_{\mathrm{out}}=534773760\ggg 15=16320.
\]
Mẫu tiếp: $y_{\mathrm{out}}=16256$. Dãy $16384,16320,16256,\ldots$ giảm khoảng $64$ mỗi bước lúc đầu --- đúng $\alpha=1-2^{-8}$. Sau khoảng $5\cdot 256$ mẫu, đầu ra tụt dưới $50$ LSB (script \texttt{verify/verify\_chuong\_03.py}).

\begin{costbox}
\textbf{Ngân sách --- DC-blocker.} Adder/sub 32~bit + dịch số học: khoảng $80$--$120$ ALM (thường vài chục đến $\sim 100$). 0 DSP (nhờ $\alpha=1-2^{-8}$). 0 M10K. Latency: 1 mẫu (pipeline Avalon-ST một nhịp). Fmax dễ $>200\,\text{MHz}$. Vì sao không cắt 16~bit sớm: tiết kiệm $\sim 30$ ALM nhưng sai chức năng (deadband) --- đắt hơn mọi LUT trên đường $F_0$.
\end{costbox}

\subsection*{3.5. Bài tập và lời giải giải thích}

\paragraph{Bài tập 3.1 (Phình $\mathcal{R}[0]$.)}
Cho $x_0=[1,-1,1,-1]$ và offset $A=2$. Tính $\sum(x_0[n]+A)^2$ và giải thích vì sao DC-blocker phải đứng trước MAC.
\paragraph{Lời giải.} $\sum(x_0+A)^2=(3)^2+(1)^2+(3)^2+(1)^2=9+1+9+1=20$, trong khi $\sum x_0^2=4$. Offset đã nhân năng lượng $\tau=0$ lên năm lần; không chặn DC thì ngưỡng voicing và đỉnh $T_0$ đều lệch.

\paragraph{Bài tập 3.2 (Đệ quy $\alpha=0{,}5$).}
Với $x=[1,1,1,1,1]$, $\alpha=0{,}5$, $y[-1]=0$, lập bảng $y[0..4]$. Nhận xét giới hạn khi $n\to\infty$.
\paragraph{Lời giải.} $1,\ 0{,}5,\ 0{,}25,\ 0{,}125,\ 0{,}0625$ --- cấp số nhân công bội $1/2$, tiến về $0$: DC thuần bị nuốt hết.

\paragraph{Bài tập 3.3 (Bit-exact step).}
Dùng công thức $y_{\mathrm{acc}}\mathrel{+}=x_{\Delta}-(y_{\mathrm{acc}}\ggg 8)$, $x_{\Delta}=(x\ll 15)-(x_{z1}\ll 15)$, tính $y_{\mathrm{out}}$ tại hai mẫu đầu của bước $16384$ (sau ba số 0).
\paragraph{Lời giải.} $y_{\mathrm{out}}[3]=16384$, $y_{\mathrm{out}}[4]=16320$ (đã kiểm script). Sai số một LSB ở đây đồng nghĩa twin và RTL lệch thang cố định điểm --- bắt buộc khớp trước khi nối FIR.
